\section{Resemblance as a Quantity to Be Optimised}\label{sec:resemblance}

\subsection{Formulation}

Suppose a weight $w(a,b)\ge0$ is given for every pair of distinct letters, with larger values meaning that $a$ and $b$ look more alike, and suppose the weight is symmetric. The design of a reciprocal substitution that changes every letter and maximises total resemblance is then the problem of finding a perfect matching of maximum total weight in the complete graph on 26 vertices. This is a classical problem solvable in polynomial time by the blossom algorithm \citep{edmonds1965}. If the substitution is allowed to contain longer cycles and the weight is not symmetric, the problem becomes an assignment problem with the diagonal forbidden, solved by the Hungarian method \citep{kuhn1955}. If only $k$ pairs are to be exchanged and the remaining letters left in place, the problem is a maximum-weight matching of exactly $k$ edges, which is again polynomial. Allowing fixed points without further constraint is useless, since the identity substitution has perfect resemblance, so some such constraint, whether a number of changed letters or a quota of demonstrated arbitrariness, is needed to make the optimisation meaningful. The point of stating the problem in this form is that the optimisation machinery is independent of where the weights come from. Weights may be derived from the geometry of rendered glyphs, from feature models of letter identification \citep{pelli2006}, or from empirical confusion matrices recording how often observers mistake one letter for another \citep{townsend1971,gilmore1979}, and the same matching procedure applies to each.

\subsection{A computable proxy and its results}

To make the problem concrete, a simple proxy was computed. Each capital letter was rendered as a bitmap on a $200\times200$ pixel canvas at an em size of 160 pixels, sharing a common baseline and centred by advance width, so that relative sizes and positions were preserved and no letter was rescaled to fill its box. The weight of a pair was the Jaccard overlap of the two bitmaps, the area of the intersection of the inked regions divided by the area of their union. This measure runs from 0 for letters with disjoint ink to 1 for identical glyphs. Four typefaces were used: DejaVu Sans, Liberation Sans, DejaVu Serif and TeX Gyre Pagella. For each, the full $26\times26$ weight matrix was computed and a maximum-weight perfect matching found with a standard implementation of the blossom algorithm. The procedure, with its code, is reproduced in Appendix~\ref{app:code}.

\begin{table}[ht]
\centering
\footnotesize
\begin{tabular}{lrrrll}
\toprule
Font & Optimal total & Mean per pair & Greedy total & Weakest pair & Strongest pair\\
\midrule
DejaVu Sans & 6.689 & 0.515 & 6.616 & MV (0.23) & OQ (0.92)\\
Liberation Sans & 6.493 & 0.499 & 6.332 & AW (0.27) & OQ (0.85)\\
DejaVu Serif & 6.638 & 0.511 & 6.418 & MV (0.28) & OQ (0.90)\\
TeX Gyre Pagella & 6.443 & 0.496 & 6.102 & MV (0.26) & OQ (0.81)\\
\bottomrule
\end{tabular}

\caption{Maximum-weight reciprocal pairings of the 26 capitals under the pixel-overlap measure, by typeface. The greedy column takes the best remaining pair repeatedly. Overlap values are Jaccard indices.}\label{tab:resemblance}
\end{table}

The totals in Table~\ref{tab:resemblance} correspond to mean overlaps of about one half per pair in every font, which is to say that on average the two letters of a pair share about half of the area covered by their union. The greedy procedure that repeatedly takes the best remaining pair comes within roughly one to five per cent of the optimum in total weight, and the gap is largest for the Pagella serif face, so a hand design that proceeds pair by pair and settles the awkward letters last is not far from optimal under this measure but can be improved by a global computation. The strongest pair in every font is O and Q, whose overlap lies between 0.81 and 0.92, and the weakest pair is M and V in three of the four fonts and A and W in the fourth.

Table~\ref{tab:named} lists the overlaps for the exchanges that arise naturally in discussion, including B and P, U and V, C and G, E and F, M and N, O and Q, I and L, and M and W. Several results deserve comment. The exchange of C and G scores between 0.54 and 0.66 and that of E and F between 0.52 and 0.65, so these are good pairs by the measure. B and P score between 0.41 and 0.52, which is moderate. U and V score between 0.19 and 0.24, M and N between 0.20 and 0.26, M and W between 0.18 and 0.23, C and K about 0.2, and I and L between 0.05 and 0.12. If one asks which of these a designer would call visually related, U and V, M and W and perhaps I and L would be on most lists, and the measure ranks them low.

\begin{table}[ht]
\centering
\small
\begin{tabular}{lrrrr}
\toprule
Pair & DejaVu Sans & Liberation Sans & DejaVu Serif & TeX Gyre Pagella\\
\midrule
O/Q & 0.92 & 0.85 & 0.90 & 0.81\\
C/G & 0.54 & 0.62 & 0.66 & 0.58\\
E/F & 0.60 & 0.52 & 0.65 & 0.56\\
B/P & 0.41 & 0.52 & 0.47 & 0.49\\
M/N & 0.26 & 0.21 & 0.20 & 0.21\\
C/K & 0.22 & 0.21 & 0.19 & 0.19\\
M/W & 0.23 & 0.19 & 0.20 & 0.18\\
A/H & 0.22 & 0.21 & 0.22 & 0.19\\
U/V & 0.24 & 0.21 & 0.23 & 0.19\\
I/L & 0.05 & 0.05 & 0.12 & 0.09\\
\bottomrule
\end{tabular}

\caption{Pixel overlap for selected pairs, by typeface.}\label{tab:named}
\end{table}

The reason is not mysterious. Overlap of ink area is a measure of coincidence, and it is high when two glyphs occupy the same region of the page, as O and Q do. It is not a measure of structural similarity. U and V share a topology, an open curve with two ascending ends, but the curved bottom of one and the pointed bottom of the other leave little shared area, and M and W are related by a reflection in a horizontal axis that coincidence of area does not detect. The measure is therefore a poor stand-in for the perceptual relation that a reader uses, and its results should be read as a demonstration of the procedure and a warning, not as a recommendation. A better weight could be derived from skeletons, from feature sets, or from human confusion data, and an experimental programme of the kind described in Section~\ref{sec:empirical} would supply the last of these directly.

\subsection{What the pairings show}

Three features of the computed pairings survive the weakness of the measure, because they follow from the structure of the optimisation and not from the particular weights. First, the optimal pairing depends strongly on the typeface. Only six of the thirteen pairs are shared by all four fonts, namely A and W, C and G, D and U, H and N, M and V, and O and Q, and a further three, B and S, I and J, and T and Y, appear in three of the four. Appendix~\ref{app:pairings} gives the full pairings. Since the shapes of letters are properties of a design and not of an alphabet, a substitution chosen for one typeface need not be good in another, and an edition printed in a different font changes the quality of its own cipher. The implication is that the resemblance of a substitution is a property of an alphabet \emph{together with} a typeface, and that a project which intends editions to be printed in whatever font a reader prefers has to treat the typeface as part of the experimental condition.

Second, the global constraint costs locally good pairs. B and P is an attractive exchange taken alone, with overlap up to 0.52, yet it appears in none of the four optimal pairings, where B is matched to S in three fonts and to R in the fourth. The reason is that a perfect matching must pair every letter, so a pair that is good in isolation can be displaced if its members are needed to resolve the pairing of difficult letters elsewhere. This is the situation that the earlier remark about letters such as A, H, R, W, X and Z anticipates, since these are the letters with few close neighbours, and the computed pairings agree: A is paired with W in all four fonts, at overlaps between 0.27 and 0.33 that are among the lowest in each pairing. If instead only some letters are changed, the matching is of exactly $k$ edges and pairs such as B and P can be retained. The number of letters changed is a design variable, and the counts in Table~\ref{tab:counting} show that the space of choices over which letters to change is large at every size.

Third, the choice between changing all letters and changing a subset is not merely a matter of optimisation but of purpose. A subset substitution is easier to read and demonstrates arbitrariness less thoroughly, while a full substitution is harder to acquire and makes a stronger point about the conventionality of the shapes. The empirical programme in Section~\ref{sec:empirical} treats the number of changed letters as one of its manipulated factors.
